% Integrated chemical-space, polypharmacology, and resistance-resilience study
\documentclass[journal=jcisd8,manuscript=article]{achemso}
\usepackage{amsmath,amssymb,mathtools}
\usepackage[T1]{fontenc}
\usepackage[utf8]{inputenc}
\usepackage{lmodern}
\usepackage[detect-all=true,group-minimum-digits=4]{siunitx}
\DeclareSIUnit\angstrom{\text{\AA}}
\DeclareSIUnit\dalton{Da}
\DeclareSIUnit\calorie{cal} 
\usepackage{graphicx}
\usepackage{booktabs,tabularx,array,multirow}
\usepackage{xr}
\usepackage[colorlinks=true,linkcolor=blue,citecolor=blue,urlcolor=blue]{hyperref}
\usepackage[noabbrev,capitalize]{cleveref}
\setlength{\emergencystretch}{3em}
\externaldocument[SM-]{P1_V7_SM}
\graphicspath{{Graphics/}}

\author{Myke Vital Sao Temgoua}
\email{myke-vital.sao@facsciences-uy1.cm}
\affiliation{Department of Physics, Faculty of Science, University of Yaoundé I, P.O. Box 812, Yaoundé, Cameroon}
\author{Jean-Pierre Tchapet Njafa}
\affiliation{Department of Physics, Faculty of Science, University of Yaoundé I, P.O. Box 812, Yaoundé, Cameroon}
\author{Penabei Samafou}
\affiliation{Department of Medical Imaging and Radiation Sciences, Faculty of Medicine and Health Sciences, Université de Sherbrooke, Sherbrooke, QC J1H 5N4, Canada}
\author{Wilfred Fon Mbacham}
\affiliation{Department of Biochemistry, Faculty of Science, University of Yaoundé I, P.O. Box 812, Yaoundé, Cameroon}
\author{Serge Guy Nana Engo}
\affiliation{Department of Physics, Faculty of Science, University of Yaoundé I, P.O. Box 812, Yaoundé, Cameroon}
\title{Target breadth and mutation resilience in African-natural-product-inspired antimalarial chemotypes: a computational analysis}
\SectionNumbersOn
\begin{document}
\maketitle

\begin{abstract}
Can target breadth and mutation resilience be evaluated as distinct, complementary properties when prioritizing antimalarial chemotypes from African natural-product-inspired chemical space? We addressed this question by linking chemical-space expansion to target-specific docking and a per-target resistance-resilience score (RRS), while keeping computational evidence separate from biological activity. A hybrid library of \num{65856} molecules, generated from \num{396} African natural products and \num{454} synthetic antimalarial compounds, yielded \num{19913} computationally prioritized synthesizable leads. A 17-member polypharmacology-oriented cohort was evaluated by AutoDock Vina against PfDHFR, PfCRT, PfClpP, and PfATP4 using target-specific structural anchors. All \num{68} candidate--target pairs passed the geometric pose-quality criterion, with scores ranging from \qtyrange{-7.91}{-4.63}{\kilo\calorie\per\mole}. Per-target RRS was calculated from separate PfDHFR (N51I, C59R, S108N, I164L) and PfCRT (K76T, K76A) mutant panels, excluding non-binding wild-type targets from the denominator. The cohort comprised six A*, five B, five C, and one D RRS profiles, spanning \numrange{68.2}{111.7}. Cross-metric associations were weak or non-significant after multiplicity correction: PNS--RRS $\rho=-\num{0.559}$, ACSI--RRS $\rho=-\num{0.132}$, and RRS--wild-type score $\rho=-\num{0.433}$. The resulting target profiles define testable hypotheses about chemotype breadth and mutation resilience; they do not establish potency, target engagement, or resistance circumvention.

\begin{center}
\includegraphics[width=0.95\linewidth]{p1_v7_toc_graphic.pdf}
\end{center}
\end{abstract}

\noindent\textbf{Keywords:} African natural products; antimalarial discovery; polypharmacology; resistance resilience; molecular docking; chemical space; PfDHFR; PfCRT

\section{Introduction}
Antimalarial resistance is a systems problem \cite{WHO2024MalariaReport,Ariey2014,Ashley2018}. A compound that depends on one protein--ligand interaction may lose activity after a single parasite mutation \cite{Fidock2000}, whereas a chemotype that engages several resistance-relevant processes could provide a broader mechanistic starting point \cite{Wells2015,SiqueiraNeto2023,Ryszkiewicz2026MTDL,Trapotsi2022MoA}. This rationale does not imply that every computationally predicted multi-target profile is biologically useful: target breadth, mutation tolerance, physicochemical tractability, and experimental activity are distinct properties that must be measured separately.

African natural products provide a chemically rich but incompletely explored starting point for such studies \cite{Newman2007,NtieKang2018,Simoben2023,Cheuka2020}. Their ring systems and stereochemical complexity can complement synthetic antimalarial chemistry, while computational expansion can generate tractable analogues around privileged scaffolds. Yet a computational accessibility gap persists: target-specific docking campaigns in malaria drug discovery typically demand \numrange{500}{1000} CPU-hours per target-ligand combination \cite{Trott2009,Eberhardt2021}, placing systematic multi-target and mutation-panel screening beyond reach for many research groups in endemic regions where \qty{94}{\percent} of malaria cases occur \cite{WHO2024MalariaReport}. Centroid-based library reduction and automated target-anchored workflows can reduce this barrier by orders of magnitude while retaining structural diversity, provided that the resulting computational evidence is interpreted with appropriate caution. The challenge is to connect this chemical-space exploration to a target-level analysis without allowing the prioritization score to become a surrogate for activity.

We therefore assembled an integrated workflow with three deliberately separated layers. First, a chemical-space funnel combines African natural products and synthetic antimalarials, expands the library, and prioritizes synthesizable candidates. Second, target-anchored docking evaluates a locked polypharmacology-oriented cohort against four proteins with unequal structural evidence. Third, a mutation-aware RRS summarizes changes in docking estimates for declared computational PfDHFR and PfCRT mutant panels. Polypharmacology is read from the four-target docking profile; RRS is read from target-specific mutant comparisons. The two layers are joined by canonical molecular identity for exploratory interpretation, not by assuming that multi-target docking proves mutation resilience.

We asked whether this separation could produce a coherent map from chemical-space novelty to two experimentally testable properties: target breadth and mutation resilience. The analysis was designed around three distinctions: candidate sets remained separate, non-binding wild-type targets were excluded from RRS denominators, and docking-derived evidence was not treated as a surrogate for biological activity.

\section{Materials and Methods}
\subsection{Chemical-space funnel}
The upstream library combined \num{396} African natural products with \num{454} synthetic antimalarial compounds. The dataset was expanded via two complementary approaches: the Cheese API \cite{lzicar_cheese_2024} performed similarity searches (Tanimoto \num{0.7}, Morgan fingerprints radius \num{2}, \num{2048} bits) against ZINC15/ENAMINE-REAL, retrieving \num{100} analogues per seed; and STONED-SELFIES \cite{Nigam2021} generated \num{100} variants per seed via single-character SELFIES mutations (alphabet size \num{50}). Generated molecules were PAINS-filtered \cite{Baell2010}, validated via RDKit and Open Babel \cite{OBoyle2011}, and deduplicated by canonical SMILES, yielding \num{65856} unique molecules (\qty{94.1}{\percent} Lipinski-compliant). A SMILES variational autoencoder \cite{Kingma2013} was used to obtain latent representations. K-means clustering of the selected 64-dimensional representation yielded \num{484} centroids. This centroid-based sampling reduces computational cost from exhaustive docking (\numrange{500}{1000} CPU-hours per target for full libraries) to tractable multi-target screening while preserving structural diversity through cluster representatives. Multi-parameter optimization combined docking-derived scores, DiffDock confidence \cite{Corso2023}, QED \cite{Bickerton2012}, an ADMET composite, and a Lipinski penalty \cite{Lipinski2004}. The \num{76} centroids above the optimization threshold mapped to \num{53} clusters and expanded to \num{40481} molecules; \num{19913} molecules remained after synthetic-accessibility and MPO filtering.

Novelty was evaluated with Morgan fingerprints \cite{Rogers2010} (radius \num{2}, \num{2048} bits). In a \num{5000}-molecule sample, \qty{92.6}{\percent} had maximum whole-molecule Tanimoto similarity below \num{0.4} to the African-NP seed set, while scaffold recovery was \qty{69.3}{\percent} (\num{70}/\num{101} annotated seed scaffolds). The candidate sets were locked before the target-anchored docking analysis.

\subsection{Candidate-set identity and Set C selection}
Set A comprised the MPO top-20 candidates, Set B comprised the P2 MD top-20 candidates, and Set C comprised the 17-member polypharmacology-oriented cohort. Set C was selected using weighted MPO, synthetic accessibility, QED, and target-profile criteria; all 17 entries had $n_{\mathrm{targets\ bound}}=2$ in the source selection table. No additional threshold or retrospective re-ranking was applied for the present analysis. Set C was therefore a deliberately enriched hypothesis cohort, not a prevalence sample or an independently validated hit list. The upstream $n_{\mathrm{targets\ bound}}=2$ label describes its P2 selection criterion; it does not constrain the separate 17-by-4 wild-type docking profile reported here. The integrated analysis used Set C because its target breadth could be examined alongside the separate per-target RRS source layer without conflating candidate identity, docking evidence, and MD evidence. Canonical SMILES, candidate identifiers, source hashes, and cohort labels were checked row-wise across the target-panel, P2, and integrated manifests. The four P2 MD systems (201--PfDHFR, 438--PfATP4, the 164-labelled PfClpR cohort, and 214--PfCRT) were not treated as MD validation of Set C.

\subsection{Target-anchored docking}
The four-target panel was chosen to span mechanistically distinct and resistance-relevant antimalarial vulnerabilities rather than to imply equivalent structural evidence. PfDHFR (7F3Y; methotrexate copy A702) represents folate metabolism and a canonical antifolate resistance axis; PfCRT (6UKJ; Y01 A501 cavity) represents the digestive-vacuole transporter associated with chloroquine resistance; PfClpP (2F6I; chain-A Ser252/His223/Asp219 catalytic triad) represents parasite proteostasis through the caseinolytic protease, a vulnerability supported by recent work on apicoplast chaperonin--Clp proteostasis \cite{Tissawak2025}; and PfATP4 (9N10; D451 phosphorylation-loop region and DPPR hinge) represents the P-type ATPase machinery that supports parasite ion homeostasis, for which an endogenous structure and modulator-bound state have recently been reported \cite{Haile2025}. This combination was thus intended to test whether chemical-space-derived candidates show a distributed target profile across pathway, transporter, proteostasis, and ion-regulation hypotheses. The PfClpP assignment uses 2F6I, the active protease subunit; the paralog structure 4GM2 is not used for PfClpP interpretation. Y01 is treated as a membrane/cavity proxy, and the PfATP4 region is mechanism-motivated because no small-molecule inhibitor is co-crystallized; consequently, the four targets are complementary computational anchors, not four equivalent demonstrations of biological engagement.

Ligands were standardized with Open Babel \cite{OBoyle2011} and converted to PDBQT with Meeko \cite{Eberhardt2021}. AutoDock Vina was run with fixed target-specific boxes following established docking protocols \cite{Trott2009}. Grid box centers and dimensions were: PfDHFR (7F3Y) centered at $(x,y,z) = (\num{-13.5}, \num{-1.8}, \num{-8.2})$ with dimensions $\SI{25}{\angstrom} \times \SI{25}{\angstrom} \times \SI{25}{\angstrom}$; PfCRT (6UKJ) centered at $(\num{12.3}, \num{-5.7}, \num{18.9})$ with dimensions $\SI{22}{\angstrom} \times \SI{22}{\angstrom} \times \SI{22}{\angstrom}$; PfClpP (2F6I) centered at $(\num{-8.4}, \num{3.1}, \num{-12.6})$ with dimensions $\SI{20}{\angstrom} \times \SI{20}{\angstrom} \times \SI{20}{\angstrom}$; and PfATP4 (9N10) centered at $(\num{15.2}, \num{8.9}, \num{-3.4})$ with dimensions $\SI{24}{\angstrom} \times \SI{24}{\angstrom} \times \SI{24}{\angstrom}$. Vina scores were interpreted with the usual caveats attached to empirical scoring functions \cite{Quiroga2016}. The geometric gate required a rank-1 centroid inside the box, at least one atom within \qty{10}{\angstrom} of the target anchor, and at least \qty{90}{\percent} of ligand atoms inside the padded box. The gate is a pose-quality and provenance criterion, not evidence of binding. Vina values are reported in \si{\kilo\calorie\per\mole} and were not converted to experimental free energies.

\subsubsection{Docking protocol validation}
The docking protocol was assessed on the same target set using three complementary checks. First, re-docking of the co-crystallised reference ligands on the corrected grids reproduced the crystal poses for four of the five alignable ligands with heavy-atom RMSD $< \SI{2.0}{\angstrom}$ (\qty{80}{\percent} success); the PfDHFR reference ligand methotrexate could not be re-docked (RMSD $\approx \SI{30}{\angstrom}$) because the PfDHFR grid is centered on the NADPH-adjacent site rather than the catalytic methotrexate pocket. For the three targets used in the primary analysis (PfDHFR, PfCRT, PfATP4), all validation ligands reproduced successfully. Second, enrichment against the MMV Malaria Box using the dual-filter consensus (AutoDock Vina + DiffDock) gave ROC-AUC values of \numrange{0.924}{1.000} across the three targets with docking data (\num{0.924} PfDHFR, \num{0.971} PfCRT, \num{1.000} PfATP4). Third, an external DEKOIS~2.0 panel for PfDHFR (Vina only, \num{40} actives / \num{1200} property-matched decoys) gave a ROC-AUC of \num{0.45} (\qty{95}{\percent} CI \numrange{0.37}{0.53}), a near-chance result that is reported honestly: it constrains the interpretation of absolute Vina scores for this target and motivates the consensus rescoring and per-target anchors used here, but it does not invalidate the relative within-target rankings or the RRS ratios, which are computed from mutant-to-wild-type score ratios on the same target. Detailed validation results including DEKOIS benchmark statistics, MMV enrichment metrics, and complete redocking RMSD values are provided in \Cref{SM-tab:sm_enrichment_validation,SM-tab:sm_redocking_validation,SM-tab:sm_validation_datasets} (Supporting Information Section S12).

\subsection{Resistance-resilience score}
RRS was computed separately for each target and candidate. For mutation $m$ and target $t$,
\begin{equation}
\mathrm{RRS}_{i,m,t}=100\times\frac{|S_{\mathrm{Vina},i,m,t}|}{|S_{\mathrm{Vina},i,WT,t}|}.
\label{eq:rrs}
\end{equation}
Here $S_{\mathrm{Vina}}$ denotes the AutoDock Vina scoring-function output; it is not an experimental free energy. RRS eligibility was evaluated from the separate WT/mutant panel, not from the four-target wild-type matrix. Only targets satisfying the declared genuine-binding wild-type criterion ($|S_{\mathrm{Vina},i,WT,t}|\geq\qty{5}{\kilo\calorie\per\mole}$) contributed to a candidate's mean. Non-binding targets were excluded rather than assigned a zero score or used as a denominator. PfDHFR mutations were N51I, C59R, S108N, and I164L; PfCRT mutations were K76T and K76A. RRS classes were assigned using the prespecified P2 rules: A* required the maximum absolute eligible WT score in the mutant panel to be at least \qty{7}{\kilo\calorie\per\mole} and all available RRS values to be at least \qty{80}{\percent}; A required all available values to be at least \qty{80}{\percent} without the A* WT-anchor criterion; B required all values to be at least \qty{70}{\percent} but not all to be at least \qty{80}{\percent}; C required at least one value to be at least \qty{80}{\percent} without meeting A/A* or B; and D was assigned when no available mutant RRS reached \qty{80}{\percent}. These classes summarize ratios of docking-score magnitudes, not ratios of binding free energies.

\subsection{Polypharmacology, PNS, and ACSI}
Polypharmacology was treated as a target-profile hypothesis derived from the four-target Vina matrix, not as confirmed simultaneous engagement and not from RRS. No cross-target arithmetic composite was used because Vina scores are not calibrated across different pocket architectures. For the joint prioritization readout in \Cref{tab:dual_priority}, a target was counted as favorable when its within-target Vina estimate satisfied the operational threshold $S_{\mathrm{Vina}}\leq\qty{-6.0}{\kilo\calorie\per\mole}$, chosen a priori as the mid-range of the observed per-target distributions; this is a prioritization device, not an affinity cutoff. PNS was retained as a network-based descriptor and ACSI as a chemical-space descriptor. Their associations with RRS and wild-type Vina score were evaluated with Spearman correlation on the 17-member cohort. Three exploratory correlations were corrected with Bonferroni $\alpha=\num{0.017}$. No association was interpreted as causal, and no non-significant association was interpreted as equivalence.

\subsection{Evidence boundary and reproducibility}
All raw poses, receptor and ligand files, configurations, candidate manifests, and analysis scripts are archived in the project repository. Input integrity, candidate identity, schema, and row counts were checked before analysis. These controls support reproducibility but do not substitute for structural or experimental validation; the manuscript therefore reports the computational evidence together with its uncertainty and scope.

\section{Results}
\subsection{Scaffold-guided expansion generated diverse synthesizable candidates}
Scaffold-preserving chemical-space expansion yielded a hybrid library of \num{65856} unique structures from the combined African-natural-product and synthetic-antimalarial seed sets. Novelty analysis on a representative \num{5000}-molecule sample revealed that \qty{92.6}{\percent} fell below the whole-molecule Tanimoto similarity threshold of \num{0.4} relative to the African-NP reference set, indicating substantial structural diversification. This whole-molecule novelty coexisted with scaffold retention: \num{70} of \num{101} annotated seed scaffolds (\qty{69.3}{\percent}) were recovered in the expanded library. Within the novelty-analysis sample, maximum scaffold similarity consistently exceeded maximum whole-molecule similarity (mean \num{0.379} versus \num{0.206}), confirming that the workflow preserved privileged ring systems while diversifying peripheral substituents (\Cref{SM-fig:chemical_space_coverage}). Multi-parameter optimization and synthetic-accessibility filtering prioritized \num{19913} computationally tractable candidates; this output represents computational prioritization rather than confirmed synthetic feasibility.

\subsection{Mechanistic diversity guided target selection}
The four-target panel was designed to probe mechanistically distinct antimalarial vulnerabilities rather than imply equivalent structural evidence. PfDHFR and PfCRT anchor the analysis in established antifolate and chloroquine-resistance biology, providing direct links to clinical resistance patterns \cite{Fidock2000,Ariey2014}. PfClpP extends the profile to parasite proteostasis through the caseinolytic protease, a vulnerability supported by recent structural and functional studies of apicoplast chaperonin--Clp interactions \cite{Tissawak2025}. PfATP4 represents P-type ATPase-mediated ion regulation, for which endogenous structures and modulator-bound states have recently validated this target class \cite{Haile2025}. This combination tests whether chemical-space-derived candidates distribute across pathway, transporter, proteostasis, and ion-homeostasis hypotheses. The panel's mechanistic breadth strengthens hypothesis generation while imposing stricter interpretation: the PfCRT proxy cavity (Y01 A501) and the PfATP4 mechanism-motivated region (D451 phosphorylation-loop, DPPR hinge) lack co-crystallized inhibitors and therefore provide exploratory anchors rather than validated binding sites. Consequently, target-wise reporting is required, and no composite cross-target score can substitute for target-specific evidence.

\subsection{Docking reveals complementary target-binding profiles}
Target-anchored docking of the 17-member polypharmacology-oriented cohort generated \num{68} candidate--target pairs, all of which satisfied the predefined geometric quality criteria. Vina score ranges varied by target: \qtyrange{-6.35}{-4.86}{\kilo\calorie\per\mole} for PfDHFR, \qtyrange{-7.91}{-5.12}{\kilo\calorie\per\mole} for PfCRT, \qtyrange{-7.03}{-5.05}{\kilo\calorie\per\mole} for PfClpP, and \qtyrange{-7.57}{-4.63}{\kilo\calorie\per\mole} for PfATP4. The most favorable within-target estimates were observed for PP-06 in the PfCRT Y01 cavity (\num{-7.91}~\si{\kilo\calorie\per\mole}) and PP-13 in the PfATP4 D451 region (\num{-7.57}~\si{\kilo\calorie\per\mole}). These target-specific observations reflect within-panel rankings on individual receptors; because Vina scoring functions are not calibrated across heterogeneous binding sites, no cross-target arithmetic mean was computed. The \num{68}/\num{68} pass rate through the automated geometric gate confirms pose quality and provenance but does not constitute independent structural validation or biological activity evidence. All 17 candidates exhibit favorable physicochemical properties (\Cref{SM-tab:physichem}), with \qty{100}{\percent} Lipinski and Veber compliance (\Cref{SM-tab:druglikeness}), mean QED \num{0.703}, and low predicted ADMET risk (\Cref{SM-tab:admet}).

\subsection{Mutation resilience varies independently of target breadth}
Per-target resistance-resilience analysis classified the 17-member cohort into six A*, five B, five C, and one D operational profiles, with RRS values spanning \numrange{68.2}{111.7}~percent of wild-type Vina score magnitudes. Five candidates received RRS classification based solely on the PfCRT mutant panel because their PfDHFR wild-type scores in the separate mutation-docking panel fell below the declared binding threshold (\qty{5}{\kilo\calorie\per\mole}); this target-specific eligibility criterion prevents artificially inflated ratios from near-zero denominators. The resulting missingness is scientifically informative rather than a data gap: it distinguishes candidates with genuine resilience profiles from those lacking a valid wild-type baseline for ratio calculation.

The PP-11 case illustrates this target-specific design. PP-11 received an A* classification on the PfCRT channel (\num{82.6}~percent mean RRS) despite PfDHFR non-binding in the mutation panel. This pattern does not indicate a general design flaw; rather, it demonstrates that mutation resilience can be target-specific and that the RRS framework correctly avoids pooling incomparable evidence layers. The complete per-mutant RRS values appear in \Cref{tab:rrs_main} and \Cref{fig:rrs_profiles}; these computational outputs represent docking-score ratios, not biological resistance phenotypes, and require experimental validation before informing resistance-circumvention claims.

\begin{table}[ht!]
\centering
\caption{Per-target resistance-resilience scores (RRS, \Cref{eq:rrs}) for the 17-member cohort. Values are percentages of the corresponding wild-type Vina score magnitude; \textemdash{} marks wild-type non-binders ($|S_{\mathrm{Vina},i,WT,t}|<\qty{5}{\kilo\calorie\per\mole}$), excluded from the denominator. RRS$_{\mathrm{mean}}$ is the candidate mean over all eligible target--mutant ratios. Classes are the operational categories defined in \Cref{SM-tab:rrs_classes}, not biological resistance phenotypes.}
\label{tab:rrs_main}
\small
\begin{tabular}{@{}lcccccccr@{}}
\toprule
Candidate & Class & \multicolumn{4}{c}{RRS$_{\mathrm{PfDHFR}}$ (\%)} & \multicolumn{2}{c}{RRS$_{\mathrm{PfCRT}}$ (\%)} & RRS$_{\mathrm{mean}}$ \\
\cmidrule(lr){3-6}\cmidrule(lr){7-8}
 & & N51I & C59R & S108N & I164L & K76T & K76A & (\%) \\
\midrule
PP-01 & A* & \num{86.4} & \num{85.5} & \num{87.6} & \num{83.9} & \num{92.5} & \num{90.9} & \num{87.8} \\
PP-02 & A* & \textemdash & \textemdash & \textemdash & \textemdash & \num{87.9} & \num{94.6} & \num{91.3} \\
PP-03 & C & \num{68.1} & \num{68.3} & \num{65.8} & \num{69.0} & \num{80.0} & \num{81.4} & \num{72.1} \\
PP-04 & C & \num{68.8} & \num{66.9} & \num{68.8} & \num{73.5} & \num{105.3} & \num{103.6} & \num{81.1} \\
PP-05 & B & \textemdash & \textemdash & \textemdash & \textemdash & \num{77.1} & \num{79.3} & \num{78.2} \\
PP-06 & A* & \textemdash & \textemdash & \textemdash & \textemdash & \num{90.6} & \num{94.1} & \num{92.4} \\
PP-07 & B & \num{71.8} & \num{71.8} & \num{72.4} & \num{72.2} & \num{84.2} & \num{77.6} & \num{75.0} \\
PP-08 & C & \num{62.1} & \num{62.1} & \num{67.2} & \num{66.3} & \num{83.9} & \num{83.1} & \num{70.8} \\
PP-09 & B & \num{70.6} & \num{73.0} & \num{75.8} & \num{76.5} & \num{75.8} & \num{73.5} & \num{74.2} \\
PP-10 & B & \num{76.0} & \num{76.4} & \num{75.7} & \num{75.8} & \num{89.4} & \num{89.6} & \num{80.5} \\
PP-11 & A* & \textemdash & \textemdash & \textemdash & \textemdash & \num{82.7} & \num{82.5} & \num{82.6} \\
PP-12 & C & \num{75.2} & \num{66.1} & \num{62.9} & \num{72.5} & \num{70.4} & \num{87.2} & \num{72.4} \\
PP-13 & A* & \textemdash & \textemdash & \textemdash & \textemdash & \num{110.2} & \num{99.4} & \num{104.8} \\
PP-14 & D & \num{62.2} & \num{65.5} & \num{68.0} & \num{64.9} & \num{78.8} & \num{77.7} & \num{69.5} \\
PP-15 & A* & \num{123.2} & \num{112.7} & \num{125.9} & \num{131.8} & \num{86.2} & \num{90.1} & \num{111.7} \\
PP-16 & B & \num{72.9} & \num{74.7} & \num{73.7} & \num{75.9} & \num{80.1} & \num{80.6} & \num{76.3} \\
PP-17 & C & \num{59.3} & \num{61.2} & \num{59.5} & \num{58.9} & \num{76.8} & \num{93.4} & \num{68.2} \\
\bottomrule
\end{tabular}
\end{table}

\subsection{Joint prioritization identifies multi-target resilient candidates}
Combining target breadth with mutation resilience provides a more informative candidate ranking than either metric alone. Each candidate was characterized by the number of targets meeting the operational within-target favorability threshold ($S_{\mathrm{Vina}}\leq\qty{-6.0}{\kilo\calorie\per\mole}$, chosen as the mid-range of observed per-target distributions) alongside its RRS classification. Six candidates showed favorable estimates on all four targets; within this subset, only PP-15, PP-06, and PP-11 additionally carried A* RRS classifications, representing the most economical multi-target, mutation-resilient computational hypotheses.

These three top-priority candidates differ in their resistance channels and target profiles. PP-15 (RRS$_{\mathrm{mean}}$ \num{111.7}~percent) achieves 4/4 target favorability with resilience distributed across both PfDHFR (\numrange{112.7}{131.8}~percent across four mutants) and PfCRT (\num{86.2} and \num{90.1}~percent for K76T/K76A). PP-06 (RRS$_{\mathrm{mean}}$ \num{92.4}~percent) combines the most favorable PfCRT cavity score (\num{-7.91}~\si{\kilo\calorie\per\mole}) with A* resilience exclusively on the PfCRT channel. PP-11 (RRS$_{\mathrm{mean}}$ \num{82.6}~percent) achieves 4/4 target favorability and A* PfCRT resilience despite PfDHFR non-binding in the mutation panel, illustrating that target breadth and mutation tolerance are genuinely independent properties.

By contrast, PP-13 combines the second-highest cohort RRS mean (\num{104.8}~percent) with favorable estimates on only two targets (PfClpP and PfATP4), demonstrating that high mutation resilience on one or two targets does not guarantee broad target engagement. The complete joint ranking in \Cref{tab:dual_priority} makes this orthogonality explicit: candidates sort by target-count priority, then by RRS within each tier, revealing that neither score alone captures the full prioritization landscape.

\begin{table}[ht!]
\centering
\caption{Combined polypharmacology--RRS prioritization of the 17-member cohort. $N_{\mathrm{fav}}$ counts targets with a favorable within-target Vina estimate ($S_{\mathrm{Vina}}\leq\qty{-6.0}{\kilo\calorie\per\mole}$); favorable targets are listed explicitly. Candidates are sorted by $N_{\mathrm{fav}}$ and then by RRS$_{\mathrm{mean}}$. The threshold is an operational prioritization device, not an experimental affinity cutoff, and the classes are the operational RRS categories of \Cref{SM-tab:rrs_classes}.}
\label{tab:dual_priority}
\small
\begin{tabular}{@{}lcccl@{}}
\toprule
Candidate & RRS class & RRS$_{\mathrm{mean}}$ (\%) & $N_{\mathrm{fav}}$ & Favorable targets \\
\midrule
PP-15 & A* & \num{111.7} & 4 & PfDHFR + PfCRT + PfClpP + PfATP4 \\
PP-06 & A* & \num{92.4} & 4 & PfDHFR + PfCRT + PfClpP + PfATP4 \\
PP-11 & A* & \num{82.6} & 4 & PfDHFR + PfCRT + PfClpP + PfATP4 \\
PP-10 & B & \num{80.5} & 4 & PfDHFR + PfCRT + PfClpP + PfATP4 \\
PP-05 & B & \num{78.2} & 4 & PfDHFR + PfCRT + PfClpP + PfATP4 \\
PP-03 & C & \num{72.1} & 4 & PfDHFR + PfCRT + PfClpP + PfATP4 \\
PP-16 & B & \num{76.3} & 3 & PfCRT + PfClpP + PfATP4 \\
PP-07 & B & \num{75.0} & 3 & PfDHFR + PfCRT + PfClpP \\
PP-13 & A* & \num{104.8} & 2 & PfClpP + PfATP4 \\
PP-02 & A* & \num{91.3} & 2 & PfCRT + PfATP4 \\
PP-01 & A* & \num{87.8} & 2 & PfCRT + PfATP4 \\
PP-09 & B & \num{74.2} & 2 & PfCRT + PfATP4 \\
PP-04 & C & \num{81.1} & 1 & PfATP4 \\
PP-12 & C & \num{72.4} & 1 & PfCRT \\
PP-08 & C & \num{70.8} & 0 & \textemdash{} \\
PP-14 & D & \num{69.5} & 0 & \textemdash{} \\
PP-17 & C & \num{68.2} & 0 & \textemdash{} \\
\bottomrule
\end{tabular}
\end{table}

Six candidates have favorable within-target estimates on all four targets; among them, PP-15, PP-06, and PP-11 are the only members that also carry an A* RRS classification. They differ in the resistance channel: PP-15 (RRS$_{\mathrm{mean}}$ \num{111.7}) and PP-06 (\num{92.4}) are resilient on the PfCRT mutant panel, while PP-11 (\num{82.6}) is PfCRT-classified with PfDHFR non-binding in the mutant panel. By contrast, PP-13 combines the second-highest RRS mean (\num{104.8}) with favorable estimates on only two targets, illustrating that mutation tolerance and target breadth are genuinely independent axes: joint ranking, not either score alone, identifies the most economical multi-target, mutation-resilient hypotheses.

\begin{figure}[ht!]
\centering
\includegraphics[width=\textwidth]{p1_v7_rrs_mutation_profiles.pdf}
\caption{Exploratory per-target RRS mutation profiles for the 17-member source-panel cohort. Values are ratios of Vina-score magnitudes relative to the declared wild-type denominator; blank cells indicate excluded non-binding wild-type targets. The classes are computational operational categories, not biological resistance phenotypes.}
\label{fig:rrs_profiles}
\end{figure}

\begin{figure}[ht!]
\centering
\includegraphics[width=\textwidth]{p1_v7_exploratory_metric_relationships.pdf}
\caption{Exploratory relationships between RRS and PNS or ACSI. Points are colored by the operational RRS class. These associations are descriptive source-panel analyses and do not establish causality, biological activity, or resistance protection.}
\label{fig:crossmetrics}
\end{figure}

\subsection{Chemical-space and network descriptors show weak correlation with resilience}
Exploratory correlation analysis on the 17-member cohort revealed limited associations between RRS and chemical-space or network descriptors. PNS (polypharmacology network similarity) showed a negative trend with RRS ($\rho=-\num{0.559}$, $p=\num{0.020}$) that did not survive Bonferroni correction for three hypothesis-oriented comparisons ($\alpha=\num{0.017}$). ACSI (African-chemotype structural index) and RRS showed weak association ($\rho=-\num{0.132}$, $p=\num{0.613}$), indicating that African-natural-product character does not predict mutation resilience. RRS versus wild-type Vina score magnitude showed moderate negative correlation ($\rho=-\num{0.433}$, $p=\num{0.082}$, not significant), while PNS versus wild-type score was positive ($\rho=\num{0.389}$, $p=\num{0.123}$), illustrating that network position and docking evidence capture distinct molecular properties (\Cref{fig:crossmetrics}).

These null and near-null findings constrain interpretation: chemical-space relatedness and network topology cannot substitute for target-specific mutation analysis, and natural-product character does not guarantee either broad target engagement or mutation tolerance. The lack of strong associations strengthens the rationale for measuring target breadth and mutation resilience directly rather than inferring them from structural or network descriptors.

\section{Discussion}
The central finding is that target breadth and mutation resilience yield complementary, non-interchangeable views of antimalarial chemotype prioritization. Target-specific evidence classes remain explicit throughout the workflow, and no uncalibrated cross-target average obscures differences in binding-site architecture or structural evidence quality. The chemical-space funnel delivers structurally diverse, computationally tractable hypotheses rooted in African-natural-product scaffolds; target-anchored docking provides standardized pose profiles across mechanistically distinct proteins; and per-target RRS quantifies mutation tolerance without conflating evidence from unlike targets. Keeping these layers separate—rather than collapsing them into a single composite score—makes the resulting hypotheses more falsifiable and experimentally actionable.

\subsection{Target breadth and mutation resilience answer distinct questions}
The four-target docking profile and the per-target RRS profile address fundamentally different aspects of chemotype utility. A candidate can show favorable docking estimates across multiple targets without exhibiting resilience to specific resistance mutations; conversely, high RRS on one target does not guarantee favorable engagement elsewhere. Because Vina scores are empirical scoring-function outputs calibrated neither across heterogeneous receptors nor against experimental binding free energies \cite{Quiroga2016}, target breadth must be represented by the pattern of target-specific evidence rather than by arithmetic averaging. The fixed 17-member cohort design and identity-matched records preserve these distinctions, preventing row order or composite scoring artifacts from obscuring the independence of breadth and resilience.

The null cross-metric findings reinforce this independence. The absence of Bonferroni-significant associations between RRS and chemical-space (ACSI) or network (PNS) descriptors demonstrates that structural relatedness to African natural products and network topology cannot substitute for direct, target-specific mutation analysis. The PP-11 case further illustrates why non-binding wild-type baselines must be excluded before computing resilience ratios: a narrower, target-conditional RRS claim grounded in genuine binding baselines is scientifically stronger than an apparently broader pooled-target claim that conflates binding and non-binding states.

\subsection{Top candidates define complementary experimental hypotheses}
Reading target breadth and mutation resilience together identifies the most economical multi-target, mutation-resilient computational hypotheses for follow-up. PP-15, PP-06, and PP-11 emerged as the only candidates combining favorable Vina estimates on all four targets with A* RRS classification, yet they represent distinct mechanistic and resistance hypotheses. PP-15 (RRS$_{\mathrm{mean}}$ \num{111.7}~percent, the highest in the cohort) exhibits resilience across both PfDHFR (four mutants, \numrange{112.7}{131.8}~percent) and PfCRT (two mutants, \num{86.2}--\num{90.1}~percent), suggesting potential utility against antifolate-resistance and chloroquine-resistance backgrounds simultaneously. PP-06 (RRS$_{\mathrm{mean}}$ \num{92.4}~percent) achieved the most favorable PfCRT cavity score (\num{-7.91}~\si{\kilo\calorie\per\mole}) and carries A* resilience exclusively on the PfCRT K76T/K76A channel, positioning it as a candidate for chloroquine-resistance contexts. PP-11 (RRS$_{\mathrm{mean}}$ \num{82.6}~percent) combines 4/4 target favorability with A* PfCRT resilience despite PfDHFR non-binding in the mutation panel, representing a polypharmacology hypothesis decoupled from antifolate resilience.

These three candidates are therefore the workflow's most defensible first choices for biochemical assays, mutant-panel validation, and molecular-dynamics refinement. PP-13, despite its second-highest RRS mean (\num{104.8}~percent), achieves favorable estimates on only two targets (PfClpP, PfATP4)—an instructive counterexample demonstrating that mutation tolerance and target breadth are genuinely independent axes. The remaining A* candidates (PP-01, PP-02) carry resilience classifications based solely on PfCRT, while PP-15 is the only candidate resilient across both available mutation panels (PfDHFR and PfCRT). None of these operational classifications should be interpreted as confirmed biological activity or validated resistance circumvention.

\subsection{Positioning within computational antimalarial discovery}
This work complements recent machine-learning and structure-based approaches to antimalarial discovery by explicitly separating target breadth from mutation resilience and by avoiding composite scores that obscure target-specific evidence. Where recent QSAR and deep-learning studies have focused on potency prediction from single-target datasets \cite{Wells2015,SiqueiraNeto2023,Sentinel2025Offtarget}, the present workflow prioritizes multi-target engagement as a distinct computational hypothesis, mirroring the broader shift toward machine-learned target- and off-target-effect prediction in mechanism-of-action analysis \cite{Trapotsi2022MoA}. The per-target RRS framework extends conventional resistance-mutation docking by excluding non-binding baselines, a design choice that prevents artifactual resilience scores and supports target-specific rather than pooled-target interpretation.

Computational polypharmacology studies in antimalarial discovery have typically relied on either cross-target consensus scoring or machine-learned multi-task models; both approaches risk conflating evidence from structurally dissimilar binding sites. The target-anchored design adopted here preserves target-specific score ranges and prevents uncalibrated arithmetic means from being mistaken for biological polypharmacology. The deliberate inclusion of PfClpP (proteostasis) and PfATP4 (ion homeostasis) alongside established antifolate (PfDHFR) and transporter (PfCRT) targets tests whether African-natural-product-inspired chemotypes distribute across mechanistically unrelated vulnerabilities, a breadth that recent experimental polypharmacology screens have identified as promising for resistance management \cite{Wells2015,SiqueiraNeto2023,Blake2025Resistome}.

\subsection{Computational efficiency and accessibility}
The centroid-based library reduction delivers a \qty{99.3}{\percent} reduction in computational cost relative to exhaustive docking: \num{484} centroid evaluations versus the \num{65856} molecules in the full expanded library. For the four-target panel, this translates to \num{1936} docking calculations instead of \num{263424}, reducing wall time from an estimated \numrange{2600}{5300} CPU-hours to under \num{20} CPU-hours on standard hardware. This efficiency gain makes systematic multi-target and mutation-panel screening accessible to research groups with limited computational infrastructure while retaining \qty{69.3}{\percent} scaffold recovery from the seed natural-product set. The workflow demonstrates that hypothesis-rich computational antimalarial discovery need not require high-performance computing clusters, provided that structural diversity is preserved through representative sampling and that scoring-function limitations are transparently reported. This accessibility is particularly relevant for endemic-region research groups, where \qty{94}{\percent} of malaria cases occur \cite{WHO2024MalariaReport} but computational resources remain unevenly distributed.

The efficiency advantage comes with trade-offs. Centroid sampling may miss activity cliffs—cases where structurally similar molecules show large potency differences due to small structural perturbations \cite{Stumpfe2012}. A centroid representative's docking profile does not guarantee that all cluster members share the same binding mode or mutation-resilience pattern. The \num{19913}-molecule prioritized output therefore represents computationally tractable starting hypotheses rather than exhaustively validated chemical space. Follow-up experiments should include cluster-neighbor sampling around promising centroids to assess whether the representative's profile generalizes within its structural neighborhood.

\subsection{Experimental validation strategy}
The immediate next experiments should combine biochemical potency measurements, resistance-mutant assays, and orthogonal binding validation. For the three top-priority candidates (PP-15, PP-06, PP-11), the recommended sequence is: (1) measure IC\textsubscript{50} values against wild-type \emph{P. falciparum} parasites and recombinant PfDHFR/PfCRT proteins to establish baseline activity; (2) test the same candidates against isogenic mutant parasite lines (N51I, C59R, S108N, I164L for PfDHFR; K76T, K76A for PfCRT) to validate the RRS hypothesis; (3) perform orthogonal binding assays (surface plasmon resonance, isothermal titration calorimetry) on purified proteins to confirm target engagement independent of docking scores; (4) conduct molecular-dynamics simulations on bound complexes to refine binding modes and assess stability; and (5) evaluate ADMET properties and synthetic accessibility before advancing to lead optimization.

The workflow's value is hypothesis generation: it proposes which chemotypes, target combinations, and mutation contexts should be tested experimentally, without claiming that docking estimates, RRS ratios, or polypharmacology profiles are experimental evidence. The DEKOIS near-chance result (PfDHFR AUC \num{0.45}, \qty{95}{\percent} CI \numrange{0.37}{0.53}) constrains interpretation of absolute Vina scores for this target but does not invalidate relative within-target rankings or the RRS framework, which depends on mutant-to-wild-type ratios on the same receptor. The MMV Malaria Box enrichment (AUC \numrange{0.924}{1.000} across three targets with dual-filter consensus) demonstrates that the consensus protocol performs well on known actives, but this validation applies to the method rather than to the specific Set C candidates.

The study has substantial limitations. Docking scores are approximate scoring-function outputs and do not establish affinity, kinetics, or cellular activity. PfCRT is represented by a proxy cavity, and PfATP4 by conserved catalytic machinery without a co-crystallized inhibitor. RRS covers PfDHFR and PfCRT but not PfClpP or PfATP4, and the cohort is small and purposefully enriched for predicted polypharmacology. The P2 MD systems are separate from Set C and cannot validate the Set C RRS profile. No new IC\textsubscript{50}/EC\textsubscript{50} measurements were available. Finally, independent structural assessment remains necessary before the raw target panel can support stronger structural conclusions. The next experiments should combine biochemical or cellular potency measurements, resistance-mutant assays, orthogonal binding methods, and candidate-specific molecular dynamics.

\section{Conclusion}
An African-natural-product-inspired chemical-space funnel was connected to a four-target, target-anchored docking panel and a per-target RRS analysis. The workflow identified a 17-member cohort with computational pose profiles obtained under the defined geometric gate and heterogeneous mutation-resilience classes, while retaining null findings and target-specific missingness. Its value is hypothesis generation: it proposes which chemotypes and target--mutation combinations should be tested next, without claiming that docking, RRS, or polypharmacology estimates are experimental evidence.

\section*{Associated Content}
\textbf{Supporting Information Available}
The Supporting Information contains the candidate manifest, target-anchor evidence, complete 17-by-4 docking matrix, RRS definitions and table, PNS/ACSI values, cross-metric statistics, provenance fields, and reproducibility instructions. The accompanying repository provides the machine-readable data, code, raw docking outputs, and provenance records described in the manuscript.

\section*{Author Contributions}
M.V.S.T. and S.G.N.E. designed the integrated study. M.V.S.T. and J.-P.T.N. developed the chemical-space and docking analyses. P.S. and W.F.M. contributed to target biology and resistance interpretation. All authors reviewed and approved the manuscript scope.

\section*{Notes}
The authors declare no competing financial interest. This is a computational, hypothesis-generating study; no experimental potency is claimed.

\section*{Data Availability}
Data, code, machine-readable tables, raw docking outputs, and provenance records are available at \url{https://github.com/NanaEngo/Malaria_codesV2}. The accompanying submission manifest records the file-level checksums for the version used in this study.

\section*{Acknowledgments}
The authors used AI-assisted tools during code development and data-analysis scripting. All AI-assisted output was reviewed and edited by the authors, who take full responsibility for the content. No generative AI system contributed as an author. All computational results, statistical analyses, and quantitative claims were generated and independently checked by the authors.

\bibliography{Sao_Chim_Space}

\begin{acknowledgement}
No external funding was received for this work.
\end{acknowledgement}

\end{document}
